\section{CREATE TABLE dan Tipe Data Umum (INT, VARCHAR, DATE, TEXT, dll.)}

Tabel adalah struktur utama penyimpanan data dalam database relasional. Setiap tabel terdiri dari kolom (\textit{field}) yang masing-masing memiliki nama dan \textbf{tipe data}, serta baris (\textit{record}) yang berisi data aktual. Perintah \code{CREATE TABLE} mendefinisikan struktur tabel: nama tabel, daftar kolom beserta tipe datanya, constraint (batasan), dan opsi lain seperti primary key dan auto-increment. Pemilihan tipe data yang tepat sangat penting karena mempengaruhi performa, penyimpanan, dan integritas data \cite{mysqlDocs}.

\begin{table}[ht]
\centering
\begin{tabular}{|P{2.5cm}|P{3cm}|P{3cm}|P{3.5cm}|}
\hline
\textbf{Tipe Data} & \textbf{Kategori} & \textbf{Kapasitas} & \textbf{Contoh Penggunaan} \\
\hline
\code{INT} & Numerik & -2M s.d. +2M & ID, jumlah, umur \\
\hline
\code{BIGINT} & Numerik & Sangat besar & ID skala besar \\
\hline
\code{DECIMAL(p,s)} & Numerik & Presisi tepat & Harga, saldo \\
\hline
\code{VARCHAR(n)} & Teks & Maks. n karakter & Nama, email \\
\hline
\code{TEXT} & Teks & Maks. 65535 byte & Deskripsi, komentar \\
\hline
\code{DATE} & Tanggal & YYYY-MM-DD & Tanggal lahir \\
\hline
\code{DATETIME} & Tanggal/Waktu & YYYY-MM-DD HH:MM:SS & Timestamp log \\
\hline
\code{BOOLEAN} & Logika & 0 atau 1 & Status aktif \\
\hline
\code{ENUM} & Pilihan & Set nilai tetap & Status (draft/publish) \\
\hline
\end{tabular}
\caption{Tipe data umum di MySQL}
\end{table}

Setiap tabel sebaiknya memiliki kolom \textbf{primary key} yang mengidentifikasi setiap baris secara unik. Kolom primary key biasanya bertipe \code{INT} atau \code{BIGINT} dengan atribut \code{AUTO\_INCREMENT} sehingga nilainya dihasilkan otomatis oleh MySQL setiap kali record baru ditambahkan. Selain primary key, kolom juga dapat memiliki constraint seperti \code{NOT NULL} (wajib diisi), \code{DEFAULT} (nilai bawaan), dan \code{UNIQUE} (tidak boleh duplikat) \cite{mysqlGettingStarted}.

\begin{webcode}[language=SQL, caption={Membuat tabel dengan berbagai tipe data dan constraint}]
CREATE TABLE produk (
  id        INT AUTO_INCREMENT PRIMARY KEY,
  nama      VARCHAR(200) NOT NULL,
  harga     DECIMAL(12,2) NOT NULL DEFAULT 0,
  stok      INT NOT NULL DEFAULT 0,
  deskripsi TEXT,
  kategori  ENUM('elektronik', 'pakaian', 'makanan', 'lainnya')
              DEFAULT 'lainnya',
  aktif     BOOLEAN DEFAULT TRUE,
  dibuat    DATETIME DEFAULT CURRENT_TIMESTAMP
);

-- Melihat struktur tabel
DESCRIBE produk;
\end{webcode}

